% ============================================================
% Weekly Meeting Report — Weeks 20--21
% Topic: Data Analysis Pipeline, Stage 1 Re-run, Stage 2 Multi-seed Plan
% ============================================================

\section*{Weekly Progress Report — Weeks 20--21}
\addcontentsline{toc}{section}{Weekly Progress Report — Weeks 20--21}

\begin{tabular}{ll}
    \textbf{Date}       & \today \\
    \textbf{Student}    & Khac Duc Giang Nguyen \\
    \textbf{Supervisor} & Dr.\ Seyed Sahand Mohammadi Ziabari \\
    \textbf{Phase}      & Analysis Infrastructure \& Stage 1 Re-run (Weeks 20--21) \\
\end{tabular}

\vspace{1em}

% ------------------------------------------------------------
\subsection*{1. Overview}

This report covers the work completed during Weeks~20--21
(26~May -- 8~June~2026).
Following the completion of Stage~2 training (job~22688135,
best T2 mAP@0.5\,=\,0.4181, FM\,=\,$-$0.0153), focus shifted to
\textbf{deepening the data analysis} of Stage~1 and Stage~2
before proceeding to Stage~3. The primary motivation was
supervisor feedback: a single-seed, single-run result is
insufficient for a credible thesis; confidence intervals across
multiple seeds and a richer set of per-epoch metrics
(precision, recall, F1) are needed.

The following was completed in this period:

\begin{enumerate}
    \item \textbf{Training scripts patched} (\texttt{train\_stage1.py},
          \texttt{train\_stage2\_ddp.py}) — both now log precision,
          recall, and F1 at the optimal confidence threshold every
          epoch, in addition to mAP@0.5. The \texttt{--seed} argument
          is now a proper CLI flag. Per-epoch checkpoint saving
          (\texttt{--save-epoch-ckpts}) was added. A \texttt{device\_ids=[local\_rank]}
          bug in the DDP wrapper (using global rank instead of local GPU
          index) was fixed in both scripts.

    \item \textbf{Stage~1 re-run completed} (job~23297129,
          30~May -- 1~June~2026) — Stage~1 re-run with the patched
          script (seed\,=\,42, per-epoch checkpoints every 5 epochs)
          finished at epoch~79 via early stopping.
          \textbf{Best mAP@0.5\,=\,0.6725} (P\,=\,0.636, R\,=\,0.661,
          F1\,=\,0.648) at epoch~49, surpassing the original
          \texttt{antiuav\_rgbt14} result by $+$0.0108.
          New T1 ceiling: $\text{mAP}_{T1}^{\text{after }T1} = 0.6725$.

    \item \textbf{Multi-seed Stage~2 plan} — three independent
          Stage~2 runs (seeds 42, 123, 999) will be submitted once
          the Stage~1 re-run finishes, to produce confidence intervals
          of the form $\text{mAP@0.5} = \mu \pm \sigma$ (n\,=\,3).

    \item \textbf{Analysis pipeline written} — five post-training
          scripts (see Section~3) cover bounding-box visualisation,
          parameter drift analysis, full P-R curves, scale-stratified
          metrics, and multi-seed CI aggregation.
\end{enumerate}

% ------------------------------------------------------------
\subsection*{2. Training Script Changes}

\subsubsection*{2.1 Precision, Recall, and F1 Logging}

The \texttt{ap\_per\_class} function (YOLOMG's metric backend)
has always returned precision ($p$), recall ($r$), and F1 ($f$)
at the confidence threshold maximising $f$, alongside the AP array.
Previously these values were silently discarded:

\begin{verbatim}
# Before (both scripts)
_, _, _, _, _, ap, _ = ap_per_class(...)   # p, r, f1 thrown away

# After
_, _, p, r, f1, ap, _ = ap_per_class(...)
prec = float(np.mean(p))
rec  = float(np.mean(r))
f1_  = float(np.mean(f1))
\end{verbatim}

The updated \texttt{results.csv} columns are:

\begin{center}\small
\begin{tabular}{lll}
\hline
\textbf{Script}             & \textbf{Old columns}                        & \textbf{New columns} \\
\hline
\texttt{train\_stage1.py}   & \texttt{mAP50, mAP50-95, lr}                & + \texttt{precision, recall, f1} \\
\texttt{train\_stage2\_ddp} & \texttt{mAP50\_T2, mAP50\_T1, lr}          & + \texttt{P\_T2, R\_T2, F1\_T2,} \\
                             &                                              & \phantom{+} \texttt{P\_T1, R\_T1, F1\_T1} \\
\hline
\end{tabular}
\end{center}

\subsubsection*{2.2 Seed as CLI Argument}

Previously \texttt{SEED\,=\,42} was hardcoded. Both scripts now
accept \texttt{--seed INT} as a command-line argument, defaulting
to~42 for backward compatibility. The seed is applied via:
\begin{verbatim}
init_seeds(SEED + rank)   # each DDP rank gets a distinct seed
\end{verbatim}
This ensures that (a) different seeds produce genuinely different
training runs, and (b) within a single run, each GPU's data
augmentation is independently seeded for maximum batch diversity.

\subsubsection*{2.3 Per-epoch Checkpoint Saving}

A new \texttt{--save-epoch-ckpts} flag saves
\texttt{epoch\_000.pt}, \texttt{epoch\_005.pt}, \ldots every
\texttt{--save-epoch-interval} epochs (default 5). This enables
post-hoc per-epoch evaluation of any metric (e.g.\ computing
P-R curves at every checkpoint, or visualising detection quality
at different training stages). It is off by default to avoid
unnecessary disk usage.

\subsubsection*{2.4 DDP \texttt{device\_ids} Bug Fix}

Both scripts previously wrapped the model with:
\begin{verbatim}
DDP(model, device_ids=[rank], ...)   # WRONG for multi-node
\end{verbatim}
The \texttt{device\_ids} argument must be the \emph{local} GPU
index on the current node, not the global rank across all nodes.
For single-node 4-GPU Snellius jobs these coincide, so the bug
was silent. Fixed to:
\begin{verbatim}
gpu_id = local_rank if local_rank != -1 else rank
DDP(model, device_ids=[gpu_id], ...)   # correct
\end{verbatim}
\texttt{local\_rank} is now passed through from \texttt{\_\_main\_\_}
to the \texttt{train()} function signature in both scripts.

% ------------------------------------------------------------
\subsection*{3. Analysis Pipeline}

Five post-training analysis scripts were written and validated
(syntax-checked; all 18 internal contract assertions pass).
All scripts target \texttt{/projects/prjs2041/uav\_code/} on
Snellius and are ready to run immediately after a training job
completes.

\begin{table}[h]
\centering
\caption{Post-training analysis scripts produced in Weeks~20--21.}
\label{tab:analysis_scripts}
\small
\begin{tabular}{p{0.30\linewidth}p{0.62\linewidth}}
\hline
\textbf{Script} & \textbf{Outputs} \\
\hline
\texttt{eval\_full\_analysis.py}
  & Full P-R curve, precision/recall/F1 at optimal threshold,
    \textbf{scale-stratified mAP@0.5} (tiny/small/normal/large),
    per-sequence mAP dot plot, summary CSV. Works on RGBT or UAV410. \\[4pt]
\texttt{visualise\_detections\_rgbt.py}
  & Annotated JPEG frames from \texttt{infrared.mp4} with predicted
    boxes (green) and GT boxes (orange), assembled MP4 video,
    per-sequence P/R summary. \\[4pt]
\texttt{visualise\_detections.py}
  & Same as above but for Anti-UAV410 (flat-JPEG format). \\[4pt]
\texttt{parameter\_drift.py}
  & Per-layer relative L2 drift and cosine similarity between Stage~1
    and Stage~2 checkpoints. Tufte dot plots grouped by backbone/neck/head. \\[4pt]
\texttt{plot\_training\_analysis.py}
  & Parses \texttt{results.csv}: derives \texttt{loss\_det},
    \texttt{kd/det} ratio, FM; produces 4-panel progress figure
    and 2-panel loss decomposition figure. \\[4pt]
\texttt{plot\_multirun\_ci.py}
  & Aggregates multiple \texttt{results.csv} files (one per seed)
    into mean\,$\pm$\,std confidence-interval plots. \\
\hline
\end{tabular}
\end{table}

\subsubsection*{3.1 Scale-Stratified Evaluation}

\texttt{eval\_full\_analysis.py} computes mAP@0.5 separately
for four UAV size categories, defined by ground-truth bounding-box
area in pixels:

\begin{center}\small
\begin{tabular}{lcc}
\hline
\textbf{Category} & \textbf{Area range (px$^2$)} & \textbf{Approx.\ side length} \\
\hline
Tiny   & $< 32^2 = 1{,}024$        & $< 32$\,px \\
Small  & $32^2$--$96^2$             & 32--96\,px \\
Normal & $96^2$--$192^2$            & 96--192\,px \\
Large  & $> 192^2$                  & $> 192$\,px \\
\hline
\end{tabular}
\end{center}

\noindent This figure directly supports the thesis narrative: the three
datasets differ sharply in their size distributions
(Anti-UAV-RGBT: 68.4\% large; CST Anti-UAV: 73.6\% small),
and scale-stratified mAP quantifies exactly how that shift
affects detection performance at each stage.

% ------------------------------------------------------------
\subsection*{4. Stage 1 Re-run (job~23297129)}

\begin{table}[h]
\centering
\caption{Stage 1 re-run submission details.}
\label{tab:s1_rerun}
\begin{tabular}{ll}
\hline
\textbf{Job ID}        & 23297129 \\
\textbf{Submitted}     & 30~May~2026 \\
\textbf{Script}        & \texttt{run\_stage1\_rerun.sh} \\
\textbf{GPUs}          & 4$\times$ A100 (DDP, \texttt{--standalone}) \\
\textbf{Time limit}    & 120\,h \\
\textbf{Seed}          & 42 \\
\textbf{Epochs}        & 100 (early stopping patience\,=\,30) \\
\textbf{Batch size}    & 16 per GPU (64 effective) \\
\textbf{Epoch ckpts}   & Every 5 epochs (\texttt{--save-epoch-ckpts}) \\
\textbf{Output dir}    & \texttt{/projects/prjs2041/runs/stage1/antiuav\_rgbt15/} \\
\hline
\end{tabular}
\end{table}

\noindent The existing Stage~1 checkpoint (\texttt{antiuav\_rgbt14/},
mAP@0.5\,=\,0.6617) is not overwritten. \texttt{increment\_path}
automatically selects the next available directory name.

\subsubsection*{4.1 Results (job~23297129 — completed 1~June~2026, 13:11 CEST)}

Job~23297129 ran for \textbf{80 epochs} (0--79) before early stopping
triggered at epoch~79 (patience\,=\,30; best mAP peaked at epoch~49
and did not improve for 30 consecutive epochs). The run completed
cleanly on node \texttt{gcn43}, 4$\times$A100, in approximately
\textbf{14\,h wall time}.

\begin{table}[h]
\centering
\caption{Stage~1 re-run final results — Anti-UAV-RGBT validation set
(job~23297129, \texttt{antiuav\_rgbt15}).}
\label{tab:s1_rerun_results}
\begin{tabular}{ll}
\hline
\textbf{Metric}                     & \textbf{Value} \\
\hline
Best mAP@0.5                        & \textbf{0.6725} \\
Precision (at best epoch)           & 0.6362 \\
Recall (at best epoch)              & 0.6609 \\
F1 (at best epoch)                  & 0.6484 \\
Epoch of best checkpoint            & 49 / 100 \\
Early stopping triggered at         & epoch~79 \\
vs.\ original run (\texttt{rgbt14}) & $+$0.0108 improvement \\
\hline
\end{tabular}
\end{table}

\noindent Selected mAP@0.5 learning curve:

\begin{center}
\small
\begin{tabular}{rr|rr|rr}
\hline
\textbf{Ep.} & \textbf{mAP@0.5} &
\textbf{Ep.} & \textbf{mAP@0.5} &
\textbf{Ep.} & \textbf{mAP@0.5} \\
\hline
 1 & 0.4070 &  20 & 0.6416 &  49 & \textbf{0.6725} \\
 3 & 0.5306 &  30 & 0.6598 &  60 & 0.6640 \\
 5 & 0.6040 &  40 & 0.6680 &  75 & 0.6521 \\
10 & 0.6213 &  45 & 0.6713 &  78 & $\approx$0.649 \\
\hline
\end{tabular}
\end{center}

\noindent The new \textbf{T1 ceiling} for all subsequent stages and
Forgetting Measure computations is:
\[
    \text{mAP}_{T1}^{\text{after }T1} = \mathbf{0.6725}
    \quad \Rightarrow \quad
    \text{FM} = \text{mAP}_{T1}^{\text{after }T2} - 0.6725
\]
This replaces the previous value of 0.6617 from \texttt{antiuav\_rgbt14}.
The improvement of $+$0.0108 is meaningful: a higher T1 ceiling means
the FM denominator is stronger, making the FM values slightly more
negative for the same absolute T1 retention. The Stage~2 multi-seed
runs will use \texttt{antiuav\_rgbt15/weights/best.pt} as the teacher
and \texttt{--t1-baseline 0.6725} for all analysis scripts.

\noindent \textbf{Actual compute:} $\approx$14\,h wall time on
4$\times$A100 ($\approx$1,930\,SBU) — faster than the original
Stage~1 run because the model converged earlier (epoch~49 vs epoch~54).

\subsubsection*{4.2 Bugs Fixed in \texttt{run\_multirun\_stage2.sh} Before Submission}

Reviewing the multi-seed launcher script revealed five bugs that
would have caused all three Stage~2 jobs to fail immediately:

\begin{enumerate}
    \item \textbf{Wrong teacher checkpoint} — \texttt{STAGE1\_WEIGHTS}
          pointed to \texttt{antiuav\_rgbt14/weights/best.pt} (the
          original run). Updated to \texttt{antiuav\_rgbt15/weights/best.pt}.

    \item \textbf{Wrong SLURM partition} — \texttt{--partition=gpu}
          does not exist on Snellius; changed to \texttt{--partition=gpu\_a100}
          (confirmed working in Stage~1 and the original Stage~2).

    \item \textbf{Time limit too short} — \texttt{--time=04:00:00}
          (4\,h) is far below the ~43\,h needed. Updated to
          \texttt{--time=72:00:00} (same as job~22688135).

    \item \textbf{Wrong GPU resource syntax} — \texttt{--gpus=4} is
          not a valid Snellius directive; changed to
          \texttt{--gres=gpu:a100:4} (matches Stage~1 script).

    \item \textbf{\texttt{torchrun} bare command} — SLURM jobs do not
          inherit the user's \texttt{PATH}, so \texttt{torchrun} would
          fail with ``command not found'' (same issue as job~22637974).
          Fixed by using the absolute path:
          \texttt{/home/knguyen1/.conda/envs/uav\_master/bin/torchrun}.
          Module loading also corrected to use
          \texttt{Miniconda3/23.5.2-0} + \texttt{source activate uav\_master}.
\end{enumerate}

% ------------------------------------------------------------
\subsection*{5. Stage 2 Multi-seed Plan}

Once job~23297129 completes and the new Stage~1 \texttt{best.pt}
is available, three Stage~2 runs will be submitted in parallel:

\begin{center}\small
\begin{tabular}{lllll}
\hline
\textbf{Run name} & \textbf{Seed} & \textbf{Teacher} & \textbf{Script} & \textbf{Est.\ SBU} \\
\hline
\texttt{seed42}  & 42  & \texttt{antiuav\_rgbt15/best.pt} & \texttt{run\_multirun\_stage2.sh} & $\approx$2,200 \\
\texttt{seed123} & 123 & same                              & same                              & $\approx$2,200 \\
\texttt{seed999} & 999 & same                              & same                              & $\approx$2,200 \\
\hline
\end{tabular}
\end{center}

\noindent Each run saves per-epoch checkpoints every 5 epochs.
After all three finish, \texttt{plot\_multirun\_ci.py} aggregates
the three \texttt{results.csv} files and produces a figure of the
form:
\[
    \text{mAP@0.5 (T2)} = \mu \pm \sigma \quad (n = 3\ \text{seeds})
\]
alongside CI bands for precision, recall, F1, and FM across epochs.

\noindent \textbf{Decision point resolved:} the Stage~1 re-run
produced mAP@0.5\,=\,0.6725, which is $+$0.0108 above the original
0.6617 --- well above the $>$0.005 threshold. The original Stage~2
results (job~22688135, teacher\,=\,\texttt{rgbt14}) are therefore
treated as \textbf{preliminary only}. The three multi-seed Stage~2
runs all use \texttt{antiuav\_rgbt15/best.pt} as the teacher,
and the T1 baseline for FM is updated to \textbf{0.6725}.

\subsubsection*{5.1 Submission — 1~June~2026}

The fixed \texttt{run\_multirun\_stage2.sh} was uploaded to Snellius
and executed. \textbf{Important:} the script must be run as a regular
bash script (\texttt{bash run\_multirun\_stage2.sh}), not submitted
directly to SLURM (\texttt{sbatch run\_multirun\_stage2.sh}).
Running it via \texttt{sbatch} submits the entire launcher as a single
CPU job (job~23340219), which then submits the three GPU jobs from
within that allocation — an unusual but valid pattern on Snellius.
The correct invocation for future runs is:
\begin{verbatim}
bash /projects/prjs2041/uav_code/run_multirun_stage2.sh
\end{verbatim}

\noindent \textbf{First submission (A100) cancelled and resubmitted on H100.}
Jobs~23340230--23340232 (gpu\_a100, pending) were cancelled and
resubmitted on the \texttt{gpu\_h100} partition for faster throughput
and shorter queue wait. H100 nodes on Snellius have 4$\times$H100
GPUs (80\,GB VRAM each) vs.\ 4$\times$A100 (40\,GB), and offer
$\approx$1.5--2$\times$ higher training throughput.

\begin{center}\small
\begin{tabular}{lllll}
\hline
\textbf{Job ID} & \textbf{Name} & \textbf{Seed} & \textbf{Partition} & \textbf{Status} \\
\hline
23349180 & \texttt{s2\_seed42}  & 42  & gpu\_h100 & PENDING \\
23349181 & \texttt{s2\_seed123} & 123 & gpu\_h100 & PENDING \\
23349182 & \texttt{s2\_seed999} & 999 & gpu\_h100 & PENDING \\
\hline
\end{tabular}
\end{center}

\noindent Each job requests 4$\times$H100, 72\,h time limit.
Actual epoch time: $\approx$34.5\,min/epoch
(vs.\ $\approx$51\,min on A100 — $\approx$1.5$\times$ speedup).
Logs:
\texttt{/projects/prjs2041/logs/stage2\_seed\{42,123,999\}\_\{JOBID\}.out}

\subsubsection*{5.2 Results (2~June~2026)}

\begin{table}[h]
\centering
\caption{Stage~2 multi-seed results (teacher\,=\,\texttt{antiuav\_rgbt15},
T1 baseline\,=\,0.6725). FM\,=\,best T1 mAP during Stage~2\,$-$\,0.6725.}
\label{tab:stage2_multiseed}
\small
\begin{tabular}{lrrrrr}
\hline
\textbf{Seed} & \textbf{Best T2 mAP} & \textbf{Best T1 mAP} & \textbf{FM} &
\textbf{Stop ep.} & \textbf{Status} \\
\hline
42  & 0.4264 & 0.6658 & $-$0.0067 & 31 & Complete \\
123 & 0.4304 & 0.6647 & $-$0.0078 & 30 & Complete \\
999 & 0.4275 & 0.6647 & $-$0.0078 & 27\,* & Cancelled (hung) \\
\hline
\textbf{Mean $\pm$ std (n\,=\,3)} & \textbf{0.428 $\pm$ 0.002} & & \textbf{$-$0.007 $\pm$ 0.001} & & \\
\hline
\multicolumn{6}{l}{\small * \texttt{best.pt} saved at epoch~27; job cancelled mid epoch~29 due to validation hang.}
\end{tabular}
\end{table}

\subsubsection*{5.3 Seed~999 hang and cancellation}

Job~23349660 (seed~999) hung during the end-of-epoch validation pass
of epoch~29. The last training-batch log entry was
\texttt{[29/49][2650/3344]} at 13:00~CEST on 2~June~2026; no further
output appeared for over 75~minutes. Each epoch normally completes in
$\approx$35\,min (training\,+\,validation), so a 75-minute silence
after batch~2650 (79\% through training) indicates the job froze in
the distributed validation \texttt{all\_gather\_object} call ---
the same NCCL deadlock pattern seen in Stage~1 job~22501881 and
Stage~2 job~22653533.

The job was cancelled (\texttt{scancel 23349660}).
Seed~999's \texttt{best.pt} (epoch~27, T2~mAP\,=\,0.4275) and
partial \texttt{results.csv} (epochs~0--28) are saved and are
sufficient for analysis. All three seeds' best checkpoints are
available, giving a valid n\,=\,3 confidence interval:
\[
    \text{mAP@0.5 (T2)} = 0.428 \pm 0.002 \quad (n = 3\ \text{seeds})
\]

\subsubsection*{5.4 Key observations}

\begin{enumerate}
    \item \textbf{T2 mAP improved vs.\ original run.}
          All three seeds achieve 0.4264--0.4304 vs.\ 0.4181 in
          job~22688135 ($+$0.008--0.012). The stronger teacher
          (\texttt{rgbt15}, mAP\,=\,0.6725 vs.\ 0.6617) provides
          richer KD guidance, raising the T2 performance ceiling.

    \item \textbf{Less forgetting than the original run.}
          FM\,$\approx$\,$-$0.007 vs.\ the original $-$0.0153.
          Absolute T1 retention is also higher ($\approx$0.665 vs.\
          $\approx$0.646), confirming the stronger teacher produces
          a more knowledge-preserving student.

    \item \textbf{Tight confidence interval.}
          T2 mAP spread of $\pm$0.002 across all three seeds confirms
          the result is stable and not seed-dependent.
\end{enumerate}

\noindent \textbf{Note on FM interpretation.}
The new FM ($\approx -0.007$) is numerically smaller in magnitude than
the original ($-0.0153$), but the T1 baseline has also increased
by 0.0108. In absolute retention terms both runs maintain T1 mAP
around 0.64--0.67 during Stage~2; the difference in FM largely reflects
the stronger baseline. Both values indicate very effective forgetting
prevention by the KD loss.

% ------------------------------------------------------------
\subsection*{6. Stage 1 Analysis — First Run and Bug Fix}

\subsubsection*{6.1 First run (job~23340483 — completed 1~June~2026, 18:07 CEST)}

A Stage~1 analysis job was submitted immediately after training completed
(1$\times$A100, 13 minutes total). The job ran two scripts:
\textbf{\texttt{eval\_full\_analysis.py}} (61,999 val frames) and
\textbf{\texttt{visualise\_detections\_rgbt.py}} (5 sequences,
80 frames each).

Overall results confirmed the training metric:

\begin{center}\small
\begin{tabular}{ll}
\hline
\textbf{Metric}  & \textbf{Value} \\
\hline
mAP@0.5          & 0.6728 \\
Precision        & 0.6360 \\
Recall           & 0.6615 \\
F1               & 0.6485 \\
Sequences        & 67 evaluated \\
GT objects       & 60,548 \\
\hline
\end{tabular}
\end{center}

\subsubsection*{6.2 Bug: Scale-Stratified Results Were All Zero}

The scale-stratified output exposed a critical bug:

\begin{verbatim}
tiny     mAP@0.5 = 0.5740   ← correct
small    mAP@0.5 = 0.7044   ← correct
normal   mAP@0.5 = 0.0000   ← WRONG
large    mAP@0.5 = 0.0000   ← WRONG
\end{verbatim}

\noindent \textbf{Root cause:} \texttt{eval\_full\_analysis.py} used
COCO-standard size bins (tiny $<32$\,px, small $32$--$96$\,px,
normal $96$--$192$\,px, large $>192$\,px). These thresholds are
appropriate for ImageNet-scale object detection but completely
wrong for UAV detection: the evaluation confirmed that
\emph{all 60,548 GT boxes in Anti-UAV-RGBT have side length $<$96\,px}.
Objects never reach the ``normal'' or ``large'' COCO thresholds
because UAVs are small targets by definition.

\noindent \textbf{Fix:} Size bins updated to UAV-appropriate thresholds
based on the longest bounding-box side:

\begin{center}\small
\begin{tabular}{lll}
\hline
\textbf{Category} & \textbf{Old (COCO)} & \textbf{New (UAV)} \\
\hline
Tiny   & $<$32\,px   & $<$16\,px \\
Small  & 32--96\,px  & 16--32\,px \\
Normal & 96--192\,px & 32--64\,px \\
Large  & $>$192\,px  & $>$64\,px \\
\hline
\end{tabular}
\end{center}

\noindent The bin definitions in \texttt{eval\_full\_analysis.py},
the caption of Figure~\ref{fig:sizes}, and the surrounding prose in
the results section were all updated to reflect these UAV-specific
thresholds. The scale distribution figure was also regenerated.

\noindent \textbf{Note:} The scale distribution percentages in
Figure~\ref{fig:sizes} (e.g.\ 68.4\% ``large'' for Anti-UAV-RGBT)
must be recomputed from the actual dataset annotations using the
new 16/32/64\,px bins via \texttt{audit\_datasets.py} on Snellius.
The current figure uses placeholder values and will be updated
once the annotation audit runs.

\subsubsection*{6.3 Re-run (job~23349514 — completed 1~June~2026, 19:18 CEST)}

After two failed attempts (KeyError on \texttt{loss\_kd} and a duplicate
\texttt{\_\_main\_\_} block from a prior heredoc corruption),
\texttt{plot\_training\_analysis.py} was fixed to auto-detect Stage~1
vs Stage~2 CSV format by checking for the \texttt{loss\_kd} column.
All three steps completed successfully in \textbf{12 minutes}
on node \texttt{gcn40} (1$\times$A100).

\paragraph{Step 1 — Training curves.}
Three figures produced from \texttt{results.csv}:
\texttt{fig\_stage1\_progress.png} (mAP, loss, P/R/F1 per epoch),
\texttt{fig\_loss\_decomp\_v2.png}, \texttt{fig\_lr\_curve.png}.
Loss fell from 0.0288 (epoch~0) to 0.0059 (epoch~79);
mAP@0.5 peaked at \textbf{0.6725} at epoch~49.

\paragraph{Step 2 — Full evaluation with corrected UAV size bins.}
Overall metrics confirmed:

\begin{center}\small
\begin{tabular}{ll}
\hline
mAP@0.5    & 0.6728 \\
Precision  & 0.6360 \\
Recall     & 0.6615 \\
F1         & 0.6485 \\
\hline
\end{tabular}
\end{center}

\noindent Scale-stratified mAP@0.5 with corrected bins
(tiny $<$16\,px, small 16--32\,px, normal 32--64\,px, large $>$64\,px):

\begin{center}\small
\begin{tabular}{lrl}
\hline
\textbf{Category} & \textbf{mAP@0.5} & \textbf{Interpretation} \\
\hline
Tiny ($<$16\,px)    & 0.009 & Nearly undetectable \\
Small (16--32\,px)  & 0.579 & Moderate difficulty \\
Normal (32--64\,px) & \textbf{0.719} & Best performance \\
Large ($>$64\,px)   & 0.461 & Lower (fewer training examples) \\
\hline
\end{tabular}
\end{center}

\noindent This is the key motivating result for Stage~3:
the model trained on Anti-UAV-RGBT performs well on normal targets
(32--64\,px) but struggles on tiny objects ($<$16\,px). CST Anti-UAV
contains predominantly small and tiny targets, so catastrophic
forgetting of normal-target features is a genuine risk.

67 sequences evaluated; top performers reached mAP@0.5\,$\approx$\,0.995;
worst performers reached near 0 (highly challenging visibility conditions).

\paragraph{Step 3 — Bounding-box visualisation.}
640 annotated frames (8 sequences $\times$ 80 frames) saved to
\texttt{analysis/vis/frames/}.
Best sequence: \texttt{20190925\_131530\_1\_3} (P\,=\,1.000, R\,=\,0.588).
Three sequences showed zero detections — likely tiny/near-invisible
targets below the 0.25 confidence threshold.
MP4 assembly skipped (\texttt{ffmpeg} not available on compute nodes).

\subsubsection*{Stage 2 multi-seed (once all three seed runs finish)}

\begin{verbatim}
# Aggregate CI plots
python plot_multirun_ci.py \
    --runs-root /projects/prjs2041/runs/stage2 \
    --pattern   "seed*" \
    --t1-baseline 0.6725   # updated: antiuav_rgbt15 best mAP@0.5

# Per-run full analysis (repeat for seed123, seed999)
python eval_full_analysis.py \
    --weights  /projects/prjs2041/runs/stage2/seed42/weights/best.pt \
    --dataset  uav410 \
    --data-root $TMPDIR/Anti-UAV410
\end{verbatim}

% ------------------------------------------------------------
\subsection*{7. Compute Budget Update}

\begin{table}[h]
\centering
\caption{Estimated SBU consumption through end of Stage~2 multi-seed.}
\label{tab:budget_w2021}
\begin{tabular}{lrr}
\hline
\textbf{Item}                           & \textbf{SBU}        & \textbf{Running total} \\
\hline
\multicolumn{3}{l}{\textit{Actual (from \texttt{accinfo} on 2~June~2026)}} \\
\hline
Consumed to date                        & 63,911              & 63,911 \\
\quad of which: all runs to Week~19     & 26,660              & — \\
\quad of which: Stage~1 re-run (23297129) & $\approx$17,152   & — \\
\quad of which: Stage~2 $\times$3 H100 seeds & $\approx$20,099 & — \\
\hline
\textbf{Remaining (budget = 120,000)}   & \textbf{56,089}     & \\
\hline
\multicolumn{3}{l}{\textit{Projected remaining spend}} \\
\hline
Seed~999 completion + analysis jobs     & $\approx$2,500    & $\approx$66,411 \\
Stage~3 (est.\ 2 runs, $\approx$30 ep.) & $\approx$8,000   & $\approx$74,411 \\
Ablations + reruns (buffer)             & $\approx$10,000   & $\approx$84,411 \\
\hline
\textbf{Projected remaining after all work} & \textbf{$\approx$35,600} & \\
\hline
\end{tabular}
\end{table}

\noindent Budget is healthy at 56,089\,SBU remaining.
Stage~2 $\times$3 seeds consumed $\approx$20,099\,SBU on H100
(vs.\ $\approx$23,760 estimated) — slightly under budget due to
early stopping. Sufficient for Stage~3, ablations, and reruns
with $\approx$35,000\,SBU to spare.

% ------------------------------------------------------------
\subsection*{8. Stage 2 Full Analysis (job~23381446 — 2~June~2026)}

Stage~2 full analysis ran on 2~June~2026 (job~23381446, 1$\times$A100,
44~minutes). Three steps: multi-seed CI aggregation, per-seed training
curves, and full evaluation of each seed's \texttt{best.pt} on
Anti-UAV410 val (94,711 frames).

\subsubsection*{8.1 Corrected FM Values}

A critical discrepancy was identified between the FM values reported
in the multi-seed CI summary and those computed by the per-seed
training analysis.

\begin{itemize}
    \item \textbf{Previously reported FM} ($\approx -0.007$): computed as
    \texttt{best T1 mAP seen at any epoch} $-$ 0.6725. This is the
    \emph{optimistic} value — it picks the highest T1 mAP ever observed
    during training, which is always epoch~0 (before any gradient steps).
    \item \textbf{Correct FM} ($\approx -0.033$): computed as
    \texttt{T1 mAP at the epoch where T2 mAP peaked} $-$ 0.6725.
    This is the honest measure of T1 retention at the point the model
    is actually checkpointed.
\end{itemize}

\begin{table}[h]
\centering
\caption{Corrected Stage~2 FM values — at peak T2 epoch.}
\label{tab:s2_fm_corrected}
\small
\begin{tabular}{lrrr}
\hline
\textbf{Seed} & \textbf{Peak T2 mAP} & \textbf{Peak T2 epoch} & \textbf{FM at peak T2} \\
\hline
42  & 0.4264 & 16 & $-$0.0369 \\
123 & 0.4304 & 15 & $-$0.0276 \\
999 & 0.4275 & 27 & $-$0.0333 \\
\hline
\textbf{Mean $\pm$ std} & \textbf{0.428 $\pm$ 0.002} & & \textbf{$-$0.033 $\pm$ 0.004} \\
\hline
\end{tabular}
\end{table}

\noindent The FM of $-0.033 \pm 0.004$ replaces the previously reported
$-0.007$ in all thesis sections. The \texttt{train\_stage3.py} script
has been updated to track FM at the best T3 epoch rather than the
all-time best T1, preventing the same bias from affecting Stage~3 results.

\subsubsection*{8.2 Scale-Stratified Performance on Anti-UAV410}

\begin{table}[h]
\centering
\caption{Stage~2 scale-stratified mAP@0.5 on Anti-UAV410 val (mean across 3 seeds).}
\label{tab:s2_scale}
\small
\begin{tabular}{lrr}
\hline
\textbf{Stratum} & \textbf{mAP@0.5 (mean)} & \textbf{Interpretation} \\
\hline
Tiny ($<$16\,px)    & 0.001 & Essentially undetectable \\
Small (16--32\,px)  & 0.278 & Moderate \\
Normal (32--64\,px) & \textbf{0.681} & Model's primary strength \\
Large ($>$64\,px)   & 0.276 & Moderate \\
\hline
\end{tabular}
\end{table}

\noindent This directly motivates Stage~3: the Stage~2 model excels at
normal-sized targets (mAP = 0.68) but cannot detect tiny targets at all.
CST Anti-UAV is 73.6\% small/tiny. Stage~3 training will force adaptation
to small/tiny targets, risking catastrophic forgetting of normal-target
detection. This is the forgetting story the herding buffer is designed to prevent.

\subsubsection*{8.3 Loss Decomposition}

The KD/det loss ratio rises from 1.36$\times$ at epoch~0 to
$\approx$2.60$\times$ at final epoch, consistent across all three seeds.
This indicates the teacher-student constraint becomes progressively
dominant relative to the detection objective as training converges —
the model is increasingly constrained by the teacher rather than
driven by the Anti-UAV410 labels. This is expected behaviour and
confirms the KD mechanism is active throughout training.

% ------------------------------------------------------------
\subsection*{9. Stage 3 Launch — 2~June~2026}

\subsubsection*{9.1 Replay Buffer Construction (job~23391452)}

Two replay buffers were built from Stage~2 \texttt{seed42/best.pt}
features on Anti-UAV-RGBT val (60,620 UAV-present frames):

\begin{table}[h]
\centering
\caption{Replay buffer construction results (job~23391452).}
\label{tab:buffers}
\small
\begin{tabular}{lrrrrrr}
\hline
\textbf{Buffer} & \textbf{Tiny} & \textbf{Small} & \textbf{Normal} & \textbf{Large} & \textbf{Total} & \textbf{Size} \\
\hline
Herding (thesis)   & 75/197 & 75/14,798 & 75/42,608 & 75/3,017 & 300 & 274\,KB \\
Random (ablation)  & 75/197 & 75/14,798 & 75/42,608 & 75/3,017 & 300 & 290\,KB \\
\hline
\end{tabular}
\end{table}

\noindent Both buffers saved to
\texttt{/projects/prjs2041/runs/stage2/seed42/}. The tiny stratum
(197 frames, 0.3\% of val) is sampled at 38\% — intentional
over-representation to protect tiny-target knowledge during Stage~3.

\subsubsection*{9.2 Training Jobs Submitted}

Three Stage~3 conditions submitted in parallel on 2~June~2026:

\begin{center}\small
\begin{tabular}{llll}
\hline
\textbf{Job ID} & \textbf{Condition} & \textbf{Purpose} & \textbf{Status} \\
\hline
23392228 & Naive baseline       & Establishes FM\_naive (no protection)      & PENDING \\
23392229 & Random stratified    & Ablation: strata balance only, no herding  & PENDING \\
23392230 & Scale-Stratified Herding & Thesis contribution: FM\_herding         & PENDING \\
\hline
\end{tabular}
\end{center}

\noindent All three jobs use Stage~2 seed42 weights as the starting
point, 50 epochs, 4$\times$H100, batch size 16. Expected runtime:
$\approx$30\,h each. The three conditions run in parallel.

\noindent \textbf{Key improvements in \texttt{train\_stage3.py} vs.\ original:}
\begin{itemize}
    \item FM tracked at best T3 epoch (not all-time best T1)
    \item Stage~2 T1 baseline measured before first gradient step
    \item Two FM values reported: \texttt{fm\_abs} (vs T1 ceiling) and
          \texttt{fm\_stage3} (vs Stage~2 T1 retention)
    \item Per-stratum T1 evaluation every epoch (tiny/small/normal/large)
    \item Random stratified ablation condition added
    \item Replay batch size increased (REPLAY\_RATIO 0.1→0.25) with
          explicit loss weighting (REPLAY\_WEIGHT=4.0)
    \item Double forward-pass bug fixed
\end{itemize}

\subsubsection*{9.3 Expected Results}

The thesis claim is FM\_herding $>$ FM\_random $>$ FM\_naive (less
negative = less forgetting). The per-stratum breakdown will show
whether herding specifically protects large/normal target detection —
the strata most threatened by CST's tiny-target distribution shift.

% ------------------------------------------------------------
\subsection*{10. Compute Budget Update (2~June~2026)}

\begin{table}[h]
\centering
\caption{Updated SBU estimate after Stage~3 launch.}
\label{tab:budget_w22}
\small
\begin{tabular}{lrr}
\hline
\textbf{Item} & \textbf{SBU} & \textbf{Running total} \\
\hline
Consumed to date (2~June)       & $\approx$65,000 & 65,000 \\
Stage~3 $\times$3 H100 (est.)   & $\approx$12,000 & 77,000 \\
Analysis jobs (drift, eval)     & $\approx$1,000  & 78,000 \\
\hline
\textbf{Projected remaining}    & \textbf{$\approx$42,000} & \\
\hline
\end{tabular}
\end{table}

\noindent Budget remains healthy. The third condition (random stratified)
adds $\approx$4,000\,SBU vs.\ the original two-condition plan but is
scientifically necessary to isolate the herding contribution.

% ------------------------------------------------------------
\subsection*{11. Stage 3 First Attempt — Disk Quota Failure and Fix}

The initial Stage~3 submission (jobs 23392228, 23392229, 23392230) failed
during epoch~0 with:

\begin{verbatim}
OSError: [Errno 122] Disk quota exceeded:
  '/scratch-local/knguyen1.23392228/pymp-...'
\end{verbatim}

\noindent \textbf{Root cause:} All three jobs were scheduled on the same
physical H100 node (shared by default on Snellius). Each job independently
ran \texttt{rsync} of both CST Anti-UAV and Anti-UAV-RGBT into its own
\texttt{\$TMPDIR} partition on the shared local scratch disk. Three jobs
$\times$ two datasets $\approx$ 300+\,GB total overwhelmed the node's
scratch quota. The dataloader workers then failed to create inter-process
communication temp files.

\noindent \textbf{Fix:} Added \texttt{\#SBATCH --exclusive} to all three
Stage~3 scripts. This guarantees each job is allocated its own dedicated
node with an independent scratch filesystem, preventing scratch collision
regardless of how many jobs are submitted in parallel.

\noindent \textbf{Resubmission (2~June~2026):}

\begin{center}\small
\begin{tabular}{lll}
\hline
\textbf{Job ID} & \textbf{Condition} & \textbf{Status} \\
\hline
23393159 & Naive baseline        & PENDING \\
23393160 & Random stratified     & PENDING \\
23393161 & Herding               & PENDING \\
\hline
\end{tabular}
\end{center}

\noindent Also confirmed during the failed run: \texttt{mAP\_T1\_after\_T2\,=\,0.6404}
for the naive condition — the Stage~2 T1 baseline is measured correctly
before any Stage~3 gradient steps. Stage~2 forgetting relative to the T1
ceiling\,=\,$-$0.0321, consistent with the corrected FM from the full
analysis ($-$0.033\,$\pm$\,0.004).

% ------------------------------------------------------------
\subsection*{12. Stage 3 Second and Third Attempts — \texttt{--exclusive} Rejected}

\subsubsection*{12.1 Second attempt (jobs 23393159--23393161) — immediate cancellation}

After the disk quota fix, \texttt{--exclusive} was added to all three scripts
and jobs resubmitted. All three were assigned to nodes (gcn91, gcn94) but
immediately cancelled by SLURM with \texttt{State: CANCELLED, ExitCode: 0:0,
Reason: None}. No log files were written. The immediate cancellation with no
error output indicates SLURM rejected the jobs at the scheduler level, most
likely because \texttt{--exclusive} is not permitted or conflicts with the
\texttt{gpu\_h100} partition policy on Snellius.

\subsubsection*{12.2 Third attempt (jobs 23393459--23393461) — same result}

A third submission after verifying script line endings (\texttt{cat -A}
confirmed no CRLF) produced identical immediate cancellations.
\texttt{sacct} confirmed all three were CANCELLED with no reason.

\subsubsection*{12.3 Root cause and fix}

The \texttt{--exclusive} SBATCH directive is not compatible with the
\texttt{gpu\_h100} partition configuration on Snellius for this account.
The correct solution is to avoid the scratch quota problem entirely by
\textbf{reading datasets directly from project storage}
(\texttt{/projects/prjs2041/datasets/}) rather than copying to
\texttt{\$TMPDIR}. This eliminates the scratch disk collision regardless
of node sharing, at the cost of slightly slower I/O. All three scripts
were updated accordingly and the \texttt{--exclusive} directive removed.

\subsubsection*{12.4 Fourth attempt (jobs 23394054--23394056) — submitted 2~June~2026}

\begin{center}\small
\begin{tabular}{lll}
\hline
\textbf{Job ID} & \textbf{Condition} & \textbf{Status} \\
\hline
23394054 & Naive baseline        & Submitted \\
23394055 & Random stratified     & Submitted \\
23394056 & Herding               & Submitted \\
\hline
\end{tabular}
\end{center}

\noindent Scripts confirmed: no \texttt{--exclusive}, datasets read from
project storage, all other parameters unchanged.

% ------------------------------------------------------------
\subsection*{13. Stage 3 Naive — First Results (3~June~2026, 00:59~CEST)}

Job~23394054 (naive baseline) has been running on node \texttt{gcn77}
since approximately 21:30~CEST on 2~June. As of 00:59~CEST on
3~June ($\approx$3.5\,h elapsed), epochs~0--7 are confirmed
complete (each epoch $\approx$25\,min on shared H100 reading from
project storage).

The random stratified (23394055) and herding (23394056) jobs remain
\texttt{PENDING} — waiting for H100 nodes to become available.

\subsubsection*{13.1 Per-epoch results — naive condition}

\begin{table}[h]
\centering
\caption{Stage~3 naive condition — per-epoch results (epochs 0--7).
T2 baseline = mAP\_T1\_after\_T2 = 0.6404. T1 ceiling = 0.6725.}
\label{tab:s3_naive_early}
\small
\begin{tabular}{rrrrrrrr}
\hline
\textbf{Ep.} & \textbf{T3 mAP} & \textbf{T1 mAP} & \textbf{fm\_abs} & \textbf{fm\_stage3} & \textbf{Tiny} & \textbf{Normal} & \textbf{Large} \\
\hline
pre & — & 0.6404 & $-$0.0321 & 0.0000 & — & — & — \\
0  & 0.0421 & 0.2064 & $-$0.4661 & $-$0.4340 & 0.041 & 0.251 & 0.010 \\
1  & 0.0655 & 0.1271 & $-$0.5454 & $-$0.5133 & 0.005 & 0.151 & 0.000 \\
2  & 0.0739 & 0.0817 & $-$0.5908 & $-$0.5587 & 0.001 & 0.092 & 0.000 \\
3  & \textbf{0.0829} & 0.0680 & $-$0.6045 & $-$0.5724 & 0.001 & 0.079 & 0.000 \\
4  & 0.0655 & 0.0571 & $-$0.6154 & $-$0.5833 & 0.000 & 0.065 & 0.000 \\
5  & 0.0705 & 0.0527 & $-$0.6198 & $-$0.5877 & 0.000 & 0.064 & 0.000 \\
6  & 0.0628 & 0.0482 & $-$0.6243 & $-$0.5922 & 0.000 & 0.055 & 0.000 \\
7  & 0.0756 & 0.0393 & $-$0.6332 & $-$0.6010 & 0.000 & 0.044 & 0.000 \\
\hline
\end{tabular}
\end{table}

\subsubsection*{13.2 Observations}

\begin{enumerate}
    \item \textbf{Catastrophic forgetting confirmed — and severe.}
    T1 mAP dropped from 0.6404 (Stage~2 baseline) to 0.039 in just
    7 epochs — a loss of 94\% of T1 knowledge. fm\_abs has already
    reached $-0.633$, far exceeding the Stage~2 FM of $-0.033$.
    This is the worst-case baseline the thesis predicts for the naive
    condition.

    \item \textbf{Large target detection → 0 by epoch 2.}
    CST Anti-UAV has 0\% large targets ($>$64\,px). Without any
    replay protection, the model loses large-target detection entirely
    within 2 epochs. By epoch 1 it is essentially zero (0.0001) and
    never recovers. This is the mechanistic story of scale-shift
    forgetting.

    \item \textbf{Normal target detection collapsing.}
    Normal-stratum mAP (32--64\,px) fell from an estimated $\approx$0.68
    at Stage~2 to 0.044 by epoch~7. The model is sacrificing its
    primary strength (normal-sized UAV detection) to adapt to
    CST's tiny targets.

    \item \textbf{T3 performance rising slowly.}
    CST mAP is improving (best 0.0829 at epoch~3), confirming
    the model is learning the new domain. The stability-plasticity
    tradeoff is visible: T3 gains come at the cost of T1 collapse.

    \item \textbf{Early stopping clock.}
    Best T3 mAP\,=\,0.0829 at epoch~3. Early stopping triggers if
    T3 mAP does not improve for 15 consecutive epochs. Expected
    stopping around epoch~18--25 if T3 plateaus.
\end{enumerate}

\subsubsection*{13.3 Significance for the thesis}

These early results are strong preliminary evidence for H1: naive
fine-tuning on a small-target dataset causes severe catastrophic
forgetting, particularly of large and normal-scale features. The
per-stratum breakdown (large → 0 in 2 epochs) provides a mechanistic
explanation that goes beyond aggregate mAP. The herding and random
stratified conditions — when they start running — need to show
materially higher T1 mAP at comparable T3 epochs to confirm the
thesis claim FM\_herding $>$ FM\_random $>$ FM\_naive.

% ------------------------------------------------------------
\subsection*{14. Stage 3 Naive — Final Results (3~June~2026, $\sim$04:00~CEST)}

Job~23394054 completed cleanly at epoch~18 via early stopping
(patience\,=\,15; best T3 mAP at epoch~3, no improvement for 15
consecutive epochs). Total wall time: approximately 7.5\,h on a
shared H100 node reading from project storage.

\subsubsection*{14.1 Final headline numbers}

\begin{table}[h]
\centering
\caption{Stage~3 naive condition — final results.}
\label{tab:s3_naive_final}
\small
\begin{tabular}{ll}
\hline
\textbf{Metric} & \textbf{Value} \\
\hline
Best T3 mAP@0.5 (CST val)         & 0.0829 \\
Best T3 epoch                      & 3 / 18 (stopped) \\
mAP\_T1\_after\_T1 (T1 ceiling)   & 0.6725 \\
mAP\_T1\_after\_T2 (Stage 3 ref)  & 0.6404 \\
\textbf{FM\_naive (fm\_abs)}       & \textbf{$-$0.6045} \\
\textbf{FM\_naive (fm\_stage3)}    & \textbf{$-$0.5724} \\
T1 mAP at final epoch (18)         & 0.0168 \\
\hline
\end{tabular}
\end{table}

\subsubsection*{14.2 Key findings}

\begin{enumerate}
    \item \textbf{Catastrophic forgetting confirmed and severe.}
    T1 mAP collapsed from 0.6404 (Stage~2 retention) to 0.017 by
    epoch~18 — a 97\% loss of T1 knowledge. FM\_naive\,=\,$-$0.6045
    (vs T1 ceiling) and $-$0.5724 (vs Stage~2 retention). This far
    exceeds the Stage~2 FM of $-$0.033, confirming that the scale
    shift in Stage~3 causes dramatically more forgetting than the
    domain shift in Stage~2.

    \item \textbf{Large target detection destroyed in 2 epochs.}
    Large-stratum mAP dropped from 0.010 at epoch~0 to 0.000 by
    epoch~2 and never recovered. CST Anti-UAV contains 0\% large
    targets, so the model receives no gradient signal to maintain
    large-target features.

    \item \textbf{Normal target detection destroyed progressively.}
    Normal-stratum mAP fell from $\approx$0.68 (Stage~2 baseline)
    to 0.019 by epoch~18. The model's primary detection strength
    is gone.

    \item \textbf{Forgetting continues beyond T3 plateau.}
    Best T3 mAP occurred at epoch~3 (0.0829). T1 mAP continued
    declining through epoch~18 (0.017), demonstrating that forgetting
    does not stop when the new task stops improving --- a hallmark of
    catastrophic forgetting.
\end{enumerate}

\subsubsection*{14.3 Thesis significance}

FM\_naive\,=\,$-$0.6045 is the lower bound the herding condition
must beat. The per-stratum data provides mechanistic evidence:
scale-shift forgetting manifests as complete loss of large and normal
target detection — the exact features absent in CST Anti-UAV. This
directly supports H1 and establishes the severity of the problem
the herding buffer is designed to solve.

% ------------------------------------------------------------
\subsection*{15. Stage 3 Herding — Submitted (3~June~2026, $\sim$04:00~CEST)}

Immediately after naive completed, the herding condition was submitted.

\begin{center}\small
\begin{tabular}{ll}
\hline
\textbf{Weights} & \texttt{runs/stage2/seed42/weights/best.pt} \\
\textbf{Buffer} & \texttt{runs/stage2/seed42/herding\_buffer.pt} (300 exemplars) \\
\textbf{Expected runtime} & $\sim$7--8\,h (based on naive timing) \\
\hline
\end{tabular}
\end{center}

\noindent \textbf{Budget note:} accinfo at submission showed 12,383\,SBU
remaining. Billing for the naive job ($\approx$7.5\,h\,$\times$\,768\,SBU/h
$\approx$\,5,760\,SBU) may not yet be reflected due to the billing delay
observed after Stage~2. Estimated true remaining: $\approx$6,600\,SBU.
Herding expected to cost $\approx$5,760\,SBU if it runs a similar number
of epochs — budget is tight but sufficient. The random stratified condition
was cancelled to preserve budget for herding, which is the primary thesis
contribution. A budget extension request of 200,000\,SBU has been drafted
for submission to SURF after the supervisor meeting on 4~June~2026.

\subsubsection*{12.5 Revised compute budget}

The first failed batch (23392228--23392230) ran for approximately 1.5--2\,h
per job before crashing (rsync + epoch~0 training), consuming an estimated
3,000--4,000\,SBU in wasted compute. The subsequently cancelled batches
consumed negligible SBU.

\begin{table}[h]
\centering
\caption{Revised SBU estimate after Stage~3 rerun overhead.}
\label{tab:budget_rerun}
\small
\begin{tabular}{lrr}
\hline
\textbf{Item} & \textbf{SBU} & \textbf{Running total} \\
\hline
Confirmed consumed (accinfo)      & 63,911  & 63,911 \\
Failed Stage~3 first batch (est.) & $\approx$3,500 & $\approx$67,400 \\
Stage~3 $\times$3 conditions (est.) & $\approx$12,000 & $\approx$79,400 \\
Remaining analysis jobs           & $\approx$1,000  & $\approx$80,400 \\
\hline
\textbf{Projected remaining}      & \textbf{$\approx$39,600} & \\
\hline
\end{tabular}
\end{table}

\noindent Budget remains sufficient. $\approx$40,000\,SBU spare provides
margin for analysis reruns, any additional experiments requested by the
supervisor on 5~June, and the buffer week before the June~26 deadline.

% ------------------------------------------------------------
\subsection*{16. Stage 3 Herding — Resubmission (3~June~2026)}

Three successive crashes were fixed before the final submission:

\begin{enumerate}
    \item \textbf{Job~23425651} — PyTorch autograd in-place conflict (DDP
    AllReduce between two backward passes). Fixed with \texttt{model.no\_sync()}
    on the first backward.
    \item \textbf{Job~23442877} — \texttt{TypeError} in logging f-string
    (\texttt{Tensor} formatted with \texttt{:.4f} without \texttt{.item()}).
    \item \textbf{Job~23443056} — Silent data bug: \texttt{ReplayDataset}
    was looking for \texttt{meta['img\_path']} which does not exist in
    \texttt{AntiUAVRGBTDataset} metas. Every replay image silently fell back
    to a zero (black) tensor, so the replay loss never decreased and provided
    no useful gradient — T1 collapsed as fast as naive. Fixed by reconstructing
    frames from \texttt{meta['seq\_dir']} and \texttt{meta['frame']} via
    \texttt{cv2.VideoCapture}, matching the original dataset's loading logic.
\end{enumerate}

\noindent Resubmitted as \textbf{job~23449442} — running on \texttt{gcn159}
(gpu\_h100). This is the final attempt within the current budget.

\subsubsection*{16.1 Early results — job~23443056 (epochs 0--1)}

\begin{table}[h]
\centering
\caption{Stage~3 herding — early per-epoch results (job~23443056).
Compare naive: T1 mAP\,=\,0.2064 at epoch~0, 0.1271 at epoch~1.}
\label{tab:s3_herding_early}
\small
\begin{tabular}{rrrrr}
\hline
\textbf{Ep.} & \textbf{T3 mAP} & \textbf{T1 mAP} & \textbf{fm\_abs} & \textbf{fm\_stage3} \\
\hline
0 & 0.0003 & 0.0411 & $-$0.6314 & $-$0.5992 \\
1 & 0.0002 & 0.0305 & $-$0.6420 & $-$0.6098 \\
\hline
\end{tabular}
\end{table}

\noindent \textbf{Warning: herding is performing worse than naive.}
At epoch~0 the naive condition retained T1 mAP\,=\,0.2064; herding
retains only 0.0411. This is the opposite of the expected direction.
A likely cause is the replay weight (\texttt{REPLAY\_WEIGHT\,=\,4.0})
being too large: with a replay loss of $\approx$0.148, the effective
replay gradient magnitude is $4\times0.148\approx0.59$, versus the
CST gradient of $\approx$0.096 — making replay $\approx$6$\times$
larger and potentially destabilising the model. Results from later
epochs (job still running) will confirm whether this corrects or worsens.

% ------------------------------------------------------------
\subsection*{17. Compute Budget (3~June~2026)}

\begin{center}\small
\begin{tabular}{ll}
\hline
Initial budget    & 120,000\,SBU \\
Used to date      & $\approx$113,140\,SBU \\
\textbf{Remaining}& \textbf{6,860\,SBU} \\
Expiry            & 2026-12-31 \\
\hline
\end{tabular}
\end{center}

\noindent Budget is critical. The three failed herding attempts
(23425651, 23442877, 23443056) consumed $\approx$5,500\,SBU in wasted
compute. At $\approx$5,760\,SBU per full herding run, the remaining
6,860\,SBU covers \textbf{exactly one more attempt} with $\approx$1,100\,SBU
to spare --- no margin for further crashes. The budget extension request
of 200,000\,SBU must be submitted to SURF immediately.

